\section{DISCUSIÓN}

La vigilancia genómica de virus en aguas residuales adquiere valor cuando convierte una señal molecular en información epidemiológica de utilidad. Los resultados de este trabajo confirman que el esquema aplicado en la PTAR Quitumbe recupera material viral de los tres agentes patógenos seleccionados con un rendimiento de caracterización condicionado por la carga viral y por la conservación de las muestras. Esta sección interpreta cada conjunto de resultados frente a la literatura revisada y al estudio antecedente local de SARS-CoV-2 \cite{MejiaCalle_2024}.

\subsection{Detección molecular en una matriz ambiental compleja}

La detección de los tres virus en distintas etapas de la investigación confirma que el muestreo pasivo, así como el flujo de concentración y extracción, recuperan material genético viral del influente de la planta. Este resultado refleja el principio establecido desde los primeros reportes de SARS-CoV-2 en aguas residuales, donde la detección de ARN viral acompañó la circulación comunitaria aun con vigilancia clínica limitada \cite{Ahmed_2020,Medema_2020}. La capacidad de detectar SARS-CoV-2, un enterovirus no polio y RSV con un mismo procedimiento respalda la viabilidad operacional de un esquema multipatógeno en un único sitio centinela, una ventaja descrita para la vigilancia basada en aguas residuales en contextos con recursos limitados \cite{Parkins_2023,Bubba_2023}.

\subsection{Carga viral, valores de Cq y rendimiento de la secuenciación}

Las variaciones observadas en los valores de Cq en las muestras explican el rendimiento desigual de la caracterización genómica. Los análisis con linaje que se logró determinar, concentraron los valores de Cq más bajos del gen N1, mientras que las muestras con Cq tardíos generaron consenso con un alto porcentaje de nucleótidos indeterminados e incluso, no se secuenciaron. Esta descripción coincide con lo reportado previamente en la PTAR Quitumbe, donde las muestras con Cq tardíos fueron menos eficientes para la secuenciación \cite{MejiaCalle_2024}, y con la observación general de que las cargas virales bajas y el ARN fragmentado limitan la recuperación genómica en aguas residuales \cite{Parkins_2023,Child_2023}. Los valores de Cq tardíos no invalidan la detección, sino que confirman la presencia del virus y permiten seguir las tendencias de circulación de este. Mediante el valor de Cq se determina el umbral práctico a partir del cual la secuenciación deja de ser informativa \cite{Ahmed_2020}. La amplificación de un solo gen diana (N1 o N2) refuerza esta interpretación, dado que la señal cercana al límite de detección se vuelve intermitente, ciertas veces se detecta un gen, otras veces el otro, y no necesariamente ambos.

\subsection{Linajes de SARS-CoV-2 y contexto epidemiológico}

Los linajes de SARS-Cov-2 identificados entre las semanas 3 y 9 de 2025 corresponde por completo a descendientes de JN.1. Los linajes JN.1, LZ.5 y LU.2.1.1 pertenecen al clado 24A, definido por JN.1, y LF.7.1.3 corresponde al clado 24H, también derivado de JN.1 \cite{Islam_2025}. Este patrón es coherente con el panorama global de comienzos de 2025, cuando JN.1 y sus descendientes dominaron la circulación y subvariantes como LP.8.1 ganaban frecuencia en las Américas \cite{Islam_2025}. La detección de JN.1 en aguas residuales de Quito reproduce el hallazgo del antecedente local, que reportó este linaje antes de su identificación clínica \cite{MejiaCalle_2024}, y confirma la capacidad de la secuenciación de aguas residuales para identificar variantes regionalmente prevalentes cuando la cobertura es suficiente \cite{Crits_Christoph_2021}.

La interpretación de estos linajes virales frente a datos clínicos locales se ve limitada por la disponibilidad de información. La vigilancia genómica de SARS-CoV-2 en Ecuador presenta una tasa de secuenciación baja respecto al total de casos y se integra principalmente mediante esfuerzos regionales, en los que JN.1 figura como linaje dominante durante 2024, pero sin detalle específico para el país en 2025 \cite{Bruno_2025}. Esta falta de datos clínicos locales dificulta la comparación directa entre la señal ambiental y la circulación documentada en pacientes, además de reforzar el valor de la vigilancia en aguas residuales como fuente complementaria y eficiente de información genómica en un contexto de secuenciación clínica limitada \cite{Bruno_2025,Parkins_2023}.

\subsection{\textit{Enterovirus}}

La detección de Coxsackievirus A en dos muestras prospectivas indica que el flujo de PCR anidada recuperó y caracterizó material de enterovirus desde la matriz ambiental. La caracterización mostró que los amplicones correspondieron a la región VP4 y no a la región VP1 que se esperaba del virus, y la identificación por BLAST, herramienta bioinformática que permite comparar un fragmento de ADN, ARN o proteína con millones de secuencias almacenadas en bases de datos, confirmó que las secuencias analizadas pertenecían a Coxsackievirus A; este desvío revela una especificidad de los cebadores de VP1 menor que la esperada en aguas residuales, un aspecto que se deberá optimizar en futuras aplicaciones del protocolo. Este hallazgo respalda el uso de las aguas residuales para la vigilancia de enterovirus no polio, una aplicación establecida de la vigilancia ambiental \cite{Bubba_2023}, y muestra que el esquema detecta circulación de enterovirus que la vigilancia clínica de parálisis flácida aguda no captura \cite{Asghar_2014}.

La ausencia de poliovirus en las 22 muestras analizadas es consistente con el contexto epidemiológico esperado. Ecuador utiliza una vacuna inactivada de este agente infeccioso y no registra circulación documentada de poliovirus, condición en la que la detección ambiental del virus resulta infrecuente y exige métodos sensibles, como se observó en países que también emplean vacuna inactivada \cite{Zurbriggen_2008,Asghar_2014}. La utilidad de la secuenciación de la región VP1 para distinguir enterovirus y poliovirus relacionados con la vacuna, fue demostrada en episodios de detección ambiental en alcantarillado urbano \cite{Klapsa_2022} y sostiene la pertinencia de mantener este blanco en la vigilancia local aun cuando no se detecte poliovirus.

\subsection{RSV}

La detección de señal de RSV en cuatro muestras sin asignación de subgrupo reproduce el comportamiento esperado para este virus en aguas residuales, dado que, la recuperación de genomas completos de RSV requiere cargas virales elevadas, con cobertura alta en muestras que superan las 10\,000 copias/mL \cite{Le_2025}, umbral que las muestras ambientales de baja carga, tal como aguas servidas, difícilmente alcanzan. La evidencia que vincula la detección de ARN de RSV en aguas residuales con la circulación de este virus en la comunidad proviene principalmente de estudios basados en técnicas de cuantificación, más que de estudios que emplean secuenciación genómica \cite{Hughes_2022}, lo que coincide con el resultado de este trabajo, donde la señal fue detectable pero insuficiente para la caracterización. La distinción entre RSV-A y RSV-B y la selección de referencias adecuadas, necesarias para la asignación de subgrupo, dependen de una cobertura genómica que no se alcanzó en estas muestras \cite{Le_2025}. En efecto, el análisis con EPI2ME no halló alineamiento con la referencia de RSV disponible, lo que impidió continuar con el análisis bioinformático y por lo tanto, con cualquier asignación de subgrupo.

\subsection{Conservación de muestras retrospectivas como factor crítico}

El análisis de las muestras retrospectivas constituye un hallazgo metodológico relevante. Las 8 muestras conservadas en DNA/RNA Shield 2X y en PBS 1X no rindieron secuencia de poliovirus ni de RSV, a diferencia de las muestras prospectivas frescas, que concentraron la totalidad de las secuenciaciones. Esto sugiere una posible degradación del ARN o carga insuficiente en el almacenamiento, y además limita el uso de estos esquemas de conservación para el análisis genómico diferido en este tipo de matriz. El resultado se relaciona con la sensibilidad del rendimiento a la integridad del ARN y a la eficiencia de recuperación señalada para el muestreo pasivo en la PTAR Quitumbe \cite{MejiaCalle_2024}, y subraya que la planificación de estudios multipatógeno debe priorizar el procesamiento temprano de las muestras frente a la conservación prolongada.

